\section{Residue orientation: a quartic in forty unrestricted integers}
\label{int:main}

The main structural improvement of this contribution is the removal of a
dedicated positivity witness.  Source~64 uses the equation $X=2u^2$ to
start its Pell argument with $X\geq0$.  We allow $X$ to have either sign,
recover $X=\pm2^R$ with $R$ odd, and select the positive branch by the
congruence $X\equiv2\pmod3$.  The congruence is obtained by affine
restrictions on two existing coordinates, so it introduces no unknown.
The preferred quartic uses the even simpler direct restriction
$X=3X_3-1$ and has $393$ expanded symbolic monomials.  We first prove the
equivalent $40$-variable version with two affine restrictions, whose monic
$X$ definition supports the best sextic and octic projections; the direct
version is specified in Subsection~\ref{int:direct}.

\begin{theorem}\label{int:theorem}
There is an explicitly specified fixed polynomial $Q_{40}^{\mathrm{aff}}$ in forty
unrestricted integer unknowns, of exact total degree four, such that
ordinary mathematics formalizable in ZFC proves
\[
  (\exists\mathbf x\in\mathbb Z^{40})\ Q_{40}^{\mathrm{aff}}(\mathbf x)=0
  \quad\Longleftrightarrow\quad
  \text{there is a finite syntactic ZFC refutation}.
\]
Consequently ZFC proves that the absence of integer zeros of $Q_{40}^{\mathrm{aff}}$
is equivalent to $\operatorname{Con}(\mathrm{ZFC})$.
Exact polynomial projections give corresponding equations in thirty-four
integers of degree six and twenty-nine integers of degree eight.
\end{theorem}

The statement imports the explicitly identified native-compiler, fixed-tuple,
K215 calibration, and Pell-rank theorems of the repository.  The new part is
the sign-uniform power-recovery argument below and its application to literal
polynomials.  It is an ordinary proof, with exact finite algebraic checks
and an independent review; it is not a new end-to-end Lean/Rocq certificate
or an exported derivation in an internally encoded ZFC calculus.

\subsection{The fixed integers and the complete quadratic system}
\label{int:literal}

Use the whole source34 \texttt{mixed-parity-unrounded-1064921} tuple.
Its ports are
\[
 \beta=2^{1064921}-1,\qquad K_0,\qquad M_C,\qquad M_F.
\]
The last port is the already shifted high mask.  All four literal integers
are supplied as lossless gzip payloads of minimal unsigned little-endian
bytes.  Their bit lengths are respectively
$1064921,1065043,1064921,1064922$; their combined uncompressed payload
is $532479$ bytes.  The included standard-library verifier reconstructs
every byte from the retained $943$-tile list and checks the original hashes.
No coefficient is selected by assuming consistency.  The relevant fixed
facts are
\begin{equation}\label{int:fixed-bounds}
 \beta\geq17,\qquad 0<M_C<\beta<M_F<2\beta,
 \qquad K_0\geq0,\qquad \beta\equiv1\pmod3.
\end{equation}

The thirty-one core unknowns are
\[
\begin{gathered}
J_3,F,\alpha,z,f,h,r,v,w_3,\tau,\eta,\zeta,y,\gamma,o,j,\\
N,E,P,c,a,X,D,C_2,L,F_2,V,g,I,B_0,\Delta.
\end{gathered}
\]
There are also three private triples $\mathbf u,\mathbf v,\mathbf e$,
for a total of forty unknowns.  In the JSON packet, $v$ is
\texttt{v\_scale}, $B_0$ is \texttt{Baux}, and the triples are
\texttt{FA\_pos1,...,3}, \texttt{FC\_pos1,...,3}, and
\texttt{eta\_zeta\_pos1,...,3}.  Every coordinate ranges over all integers.
Define only the following polynomial abbreviations:
\begin{align*}
 J_0&=3J_3+1,& W&=4(3w_3+1),&q&=\beta J_0+1,\\
 C&=q-F-\alpha,& k&=\eta+\zeta,\\
 \mathcal R&=(q-F)(N-q)-CL+(M_C+qM_F)J_0.
\end{align*}
Also set $\operatorname{Pos}(x_1,x_2,x_3)=1+x_1^2+x_2^2+x_3(x_3+1)$.
The exact main equation is $Q_{40}^{\mathrm{aff}}=\sum_{\nu=1}^{25}r_\nu^2=0$,
where the twenty-five residuals are the following quadratic polynomials:
\begin{align*}
 r_1&=L-q^2+1,                         &r_2&=N-q(L+1),\\
 r_3&=X-qW,                           &r_4&=a-E-(2v+1)N,\\
 r_5&=E(X+1)-Xa,                      &r_6&=c-\eta+(E-a)k,\\
 r_7&=P-k-E(c-\eta),                  &r_8&=\tau^2-P^2+Pk-1,\\
 r_9&=D-X-ac-\gamma(4a+3),            &r_{10}&=\Delta-a^2-4a-3,\\
 r_{11}&=C_2-c^2,                     &r_{12}&=D^2-C_2\Delta-1.
\end{align*}
\begin{align*}
 r_{13}&=k-1-hE-\mathcal R,
   &r_{14}&=(K_0+W)C-z(q-1)-F,\\
 r_{15}&=F_2-f^2,
   &r_{16}&=V+c-fo,\\
 r_{17}&=V-jc+k-1-hE,
   &r_{18}&=V^2+g-y^2,\\
 r_{19}&=I-(2r+1)C_2,
   &r_{20}&=B_0-I^2,\\
 r_{21}&=B_0-\Delta(F_2-1),
   &r_{22}&=y^2-1-B_0g,\\
 r_{23}&=F\alpha-\operatorname{Pos}(\mathbf u),
   &r_{24}&=FC-\operatorname{Pos}(\mathbf v),\\
 r_{25}&=\eta\zeta-\operatorname{Pos}(\mathbf e).
\end{align*}
No abbreviation adds an existential unknown.  Expanding these rows gives
exactly the literal packet \texttt{quartic40.json}.  Each row has degree
two.  Only $r_8$ contains $\tau$, and its coefficient of $\tau^2$ is one;
therefore $[\tau^4]Q_{40}^{\mathrm{aff}}=1$, proving exact degree four on the actual
fixed-coefficient slice.

\subsection{What changed from the parent}
\label{int:parent}

In source64's $41$-variable quartic, make the affine substitutions
\[
 J_{\rm old}=3J_3+1,\qquad
 w_{{\rm scale},\rm old}=3w_3+1,
\]
then delete the residual $X-2u^2$ and its private coordinate $u$.
The other twenty-five rows are unchanged under those substitutions.
Since $\beta\equiv1\pmod3$, the new rows imply
\begin{equation}\label{int:orientation-residue}
 q\equiv2\pmod3,\qquad W\equiv1\pmod3,
 \qquad X=qW\equiv2\pmod3.
\end{equation}
The work is to prove that this residue restriction can replace the deleted
positivity witness without introducing spurious zeros.

\subsection{Bounds before any sign is chosen}
\label{int:bounds}

At an integer zero every residual is zero.  The form
$\operatorname{Pos}$ is strictly positive on $\mathbb Z^3$ and represents
every positive integer.  To see surjectivity, apply Legendre's three-square
theorem to $4n-3$, which is $1$ modulo $4$ and cannot have the exceptional
form $4^a(8b+7)$; its representation has one odd and two even squares.

The two guards $F\alpha>0$ and $FC>0$ force $F,\alpha,C$ to have a common
nonzero sign.  Thus $q=F+\alpha+C$ has that sign and $|q|\geq3$.
Write $Q=|q|$.  Because $q=\beta J_0+1$, one has $J_0\neq0$ and
$Q\geq\beta-1\geq16$.  Set
\[
 R=((q-1)(q-F)+\alpha)(q^2-1)+(M_C+qM_F)J_0.
\]
The first two rows give $L=q^2-1$, $N=q^3$, and hence $\mathcal R=R$.
The sign-free packing estimate is
\begin{equation}\label{int:R-bounds}
       Q(Q^2-1)<R<Q^4-2.
\end{equation}
Here is a direct reminder of its proof.  If $q=Q>0$, put
$G=Q(Q-F)-C$.  Then
\[
 Q+2\leq G\leq Q^2-Q-1,
 \quad0<(M_C+QM_F)J_0<(2Q+1)(Q-1).
\]
If $q=-Q<0$, put $F_0=-F$, $C_0=-C$; now
$G=Q(Q-F_0)+C_0$ satisfies
\[
 2Q+1\leq G\leq Q^2-2,
 \quad0<(QM_F-M_C)(Q+1)/\beta<2Q(Q+1).
\]
Multiplying the displayed bounds for $G$ by $Q^2-1$ gives
\eqref{int:R-bounds} in both cases.  This uses no sign of $X,Y,k$.

Let $s=2v+1$, $i=2r+1$, and $Y=sq^3$.  Rows $r_4,r_5$ give
$a=E+Y$ and $E=XY$; thus $a=(X+1)Y$.
The nonzero integer $W=12w_3+4$ has $|W|\geq4$.  Since $s$ is odd,
writing $U=|X|$, $Z=|Y|$, and $M=|E|$ gives
\begin{equation}\label{int:UZ-bounds}
 U\geq4Q\geq64,\qquad Z\geq Q^3\geq4096,
 \qquad M=UZ\geq4Q^4>3R+1.
\end{equation}
The last guard implies that $\eta,\zeta$ have a common nonzero sign.
Put
\[
 \xi=\operatorname{sgn}X,\quad
 \sigma=\operatorname{sgn}k,\quad
 \delta=\operatorname{sgn}Y,\quad
 K=|k|,\quad \eta_0=|\eta|,\quad c_0=|c|.
\]
Then $K\geq2$, $0<\eta_0<K$, and
\begin{equation}\label{int:c-ratio}
 c_0=KZ+\delta\eta_0,\qquad
 K(Z-1)<c_0<K(Z+1),\qquad
 \operatorname{sgn}c=\sigma\delta.
\end{equation}
Set $A=a+2$, $H=4a+3$, and define the positive absolute Pell parameters
\[
 B=|A|=(U+\xi)Z+2\xi\delta,\qquad
 P_*=|2XY^2+1|=2UZ^2+\xi.
\]
Their signs in these expressions are valid because $(U-1)Z>2$.
They satisfy
\begin{equation}\label{int:BP-bounds}
 B\geq(4Q-1)Q^3-2>Q^4>R,\qquad P_*>B\geq2,\qquad c_0>K.
\end{equation}
For instance $P_*-B\geq2UZ^2-(U+1)Z-3>0$ in
\eqref{int:UZ-bounds}.  These inequalities hold in every sign branch.

\subsection{A sign-uniform Pell power-recovery argument}
\label{int:signed-power}

For $b\geq2$, use the standard Pell polynomials
\[
 \chi_b(t)+\psi_b(t)\sqrt{b^2-1}
       =(b+\sqrt{b^2-1})^t.
\]
The following proof retains all signs until the last congruence.

\paragraph{Classify the first and main norms.}
The source rows imply
\[
 \tau^2-XY^2(XY^2+1)k^2=1,\qquad
 D^2-(B^2-1)c^2=1.
\]
Consecutive-product Pell classification gives positive indices $n,p$ with
\begin{equation}\label{int:classify}
 |\tau|=\chi_{P_*}(n),\quad K=2\psi_{P_*}(n),\quad
 |D|=\chi_B(p),\quad c_0=\psi_B(p).
\end{equation}
There is no hidden positive-$X$ hypothesis: if $X<0$, put $m=UZ^2$;
the first coefficient is $m(m-1)=(m-1)m$ and its positive Pell parameter
is $2m-1=P_*$.  Since $P_*>B$ and $c_0>K$, monotonicity of $\psi$
implies $p>n$.  Otherwise
$c_0=\psi_B(p)\leq\psi_{P_*}(n)=K/2$, a contradiction.

\paragraph{Establish the ordinary-rank hypotheses.}
Modulo $M=|E|$, the parameter $P_*$ is congruent to $\xi$.
The recurrence for $\psi$ gives
$\psi_{P_*}(n)\equiv\xi^{n-1}n\pmod M$.
The index row $k=R+1+hE$ therefore gives
\begin{equation}\label{int:first-index}
 \sigma\xi^{n-1}2n\equiv R+1\pmod M.
\end{equation}
If the sign on the left is positive, $n\geq(R+1)/2$.
If it is negative, then $2n\geq M-R-1>2R$, so $n>R$.
Thus in both branches $p>n\geq(R+1)/2\geq24$.
The elementary bound $\psi_B(p)\geq(2B-1)^{p-1}$ gives
\begin{equation}\label{int:rank-hypotheses}
 c_0>B(B^2-1)^2,\qquad c_0>2p,\qquad c_0>2R.
\end{equation}
The first already follows for $p\geq6$; the last also follows from
$c_0\geq\psi_B(2)=2B>2R$.

The ordinary strong equation is
\[
 (ic^2)^2=(B^2-1)(f^2-1),\qquad i=2r+1\neq0.
\]
The inherited ordinary strong-rank lemma, now with its size hypothesis
proved in every sign branch, supplies an $m$ such that
\begin{equation}\label{int:strong-rank}
 |f|=\chi_B(m),\qquad p\mid m,\qquad c_0\mid m,\qquad m>2p.
\end{equation}
The exact source is the retained
\texttt{PELL\_RELAXED\_AUXILIARY\_PROOF.md}; its signed use and the
next step are detailed in the retained odd-scale and signed-domain notes.
Only $B,c_0,|f|,|i|$ occur in its hypotheses.

\paragraph{Recover the main rank before orienting signs.}
Let $\rho=|i|c_0^2$.  The auxiliary equations and exposed quotients give
\[
 \rho^2V^2-(\rho^2-1)y^2=1,\qquad
 V\equiv-c\pmod{|f|},\qquad V\equiv-R\pmod{c_0}.
\]
For an odd auxiliary index $\ell=2t+1$ and a sign $e_0$, Pell
classification gives $V=e_0 Q_t(\rho^2)$, where
\[
 Q_t(1-B^2)=(-1)^t\psi_B(2t+1),\qquad
 Q_t(0)=(-1)^t(2t+1).
\]
Here $\rho^2\equiv1-B^2\pmod{|f|}$ and
$\rho^2\equiv0\pmod{c_0}$.  Squaring the first quotient congruence
removes both $e_0$ and the sign of $c$.  The inherited signed Pell
step-down theorem, using $0<2p<m$, gives
$\ell\equiv\pm p\pmod m$; the second congruence gives
$\ell\equiv\pm R\pmod{c_0}$.
Since $c_0\mid m$, one obtains $R\equiv\pm p\pmod{c_0}$.
The intervals in \eqref{int:rank-hypotheses} imply
\begin{equation}\label{int:p-equals-R}
                         p=R.
\end{equation}
This is precisely the sign-insensitive auxiliary interface used by the
parent, with its prerequisites newly established also for $X<0$.

\paragraph{Unwrap the first index.}
Since $n<p=R$, the negative-sign branch of \eqref{int:first-index},
which required $n>R$, is impossible.  In the positive branch write
$2n=R+1+v_0M$, with $v_0\geq0$.  If $v_0\geq1$, then
$M>3R+1$ again implies $n>R$.  Therefore
\begin{equation}\label{int:unwrap}
 \sigma\xi^{n-1}=1,\qquad 2n=R+1,\qquad R\text{ is odd}.
\end{equation}

\paragraph{One ratio estimate handles all signs.}
Put $r_0=n-1=(R-1)/2\geq2$.  The elementary upper and lower Pell
growth bounds give
\[
 \frac{c_0}{K}
 =\frac{\psi_B(2r_0+1)}{2\psi_{P_*}(r_0+1)}
 \geq\frac12\left(\frac{(2B-1)^2}{2P_*}\right)^{r_0}.
\]
For $U\geq64$ and $Z\geq6$,
\begin{equation}\label{int:uniform-ratio}
               \frac{(2B-1)^2}{2P_*}>U-3.
\end{equation}
Indeed $B\geq(U-1)Z-2$ and $P_*\leq2UZ^2+1$, and
\begin{align*}
 &(2(U-1)Z-5)^2-(U-3)(4UZ^2+2)\\
 &\qquad=4(U+1)Z^2-20(U-1)Z+31-2U\\
 &\qquad=U(4Z^2-20Z-2)+(4Z^2+20Z+31)>0.
\end{align*}
Both final parentheses are positive for $Z\geq6$.
Together with \eqref{int:c-ratio}, this implies
\[
 Z+1>\frac{(U-3)^{r_0}}2\geq\frac{61^{r_0}}2.
\]
As $61^{r_0}/2-1>2^{2r_0+1}$ for $r_0\geq2$, we obtain
\begin{equation}\label{int:large-Y}
                            Z>2^R.
\end{equation}
The verifier checks the exact polynomial identity and positive coefficient
certificates after shifting $U=64+u$, $Z=6+z$; the displayed argument is
the universal proof rather than a finite-box test.

\paragraph{Exclude reciprocal roots and recover a signed power.}
Let $e=\operatorname{sgn}c$ and write $D=\lambda\chi_A(R)$ with
$\lambda\in\{1,-1\}$.  Since $R$ is odd,
$\psi_A(R)=\psi_B(R)$, even for $A<0$; hence $c=e\psi_A(R)$.
The polynomial Pell congruences, valid for either sign of $A,H$, are
\begin{equation}\label{int:root-congruences}
 \chi_A(R)-a\psi_A(R)\equiv2^R\pmod H,\qquad
 \chi_A(R)+a\psi_A(R)\equiv2^{-R}\pmod H.
\end{equation}
The odd modulus makes $2$ invertible.  At index one the two congruences
are $A-a=2$ and $2(A+a)\equiv1\pmod H$; the recurrence proves all indices.

If $\lambda=-e$, the row $D-ac=X+\gamma H$ gives
$X2^R\equiv-e\pmod H$.  But
\[
 1<U2^R<UZ<4(U-1)Z-4\leq|H|-1,
\]
by \eqref{int:large-Y}.  Neither sign of $X2^R$ can be congruent to
either unit modulo $|H|$.  This branch is impossible.
If $\lambda=e$, the same row gives $X\equiv e2^R\pmod H$; now
\[
 |X-e2^R|\leq U+2^R<U+Z<4(U-1)Z-3\leq|H|.
\]
This multiple of $H$ must be zero.  We have proved, before assuming
any positive sign of $X$,
\begin{equation}\label{int:signed-power-conclusion}
                      X=e2^R,\qquad R\text{ odd}.
\end{equation}

\subsection{Orientation, soundness, and canonical completeness}
\label{int:finish}

Because $R$ is odd, $2^R\equiv2\pmod3$, while
$-2^R\equiv1\pmod3$.  Combining
\eqref{int:signed-power-conclusion} with
\eqref{int:orientation-residue} forces $e=1$ and $X=2^R$.
The integer
\[
                         u=2^{(R-1)/2}
\]
therefore restores the deleted row $X=2u^2$.
Restore also $J_{\rm old}=3J_3+1$ and
$w_{{\rm scale},\rm old}=3w_3+1$.
Every parent residual now vanishes literally.  The source64 quartic41
theorem gives the finite ZFC refutation.  Its theorem is invoked only
after its square guard has been restored, so no positive-$X$ argument
has been used circularly.

Conversely, a finite refutation supplies the canonical source34 witnesses
used in quartic41 completeness.  With $d=1064921$, they satisfy
\[
 \begin{gathered}
 q=2^t,\quad t\text{ odd},\quad R\equiv3\pmod4,\quad
 X=2^R,\quad R-t\geq2\text{ even},\\
 J_{\rm old}=\frac{2^t-1}{2^d-1},\qquad
 w_{{\rm scale},\rm old}=2^{R-t-2}.
 \end{gathered}
\]
Modulo $3$ the first old coordinate is $1$, since $\beta\equiv1$
and $2^t\equiv2$.  The second is $1$ because its exponent is
nonnegative and even.  Thus
\[
 J_3=\frac{J_{\rm old}-1}{3},\qquad
 w_3=\frac{w_{{\rm scale},\rm old}-1}{3}
\]
are integers.  Retain all other coordinates and discard $u$ to obtain a
zero of $Q_{40}^{\mathrm{aff}}$.  It is the canonical completeness family that is
preserved; no claim is made that every parent zero lifts to the two
affine restrictions.  This proves Theorem~\ref{int:theorem} for $Q_{40}^{\mathrm{aff}}$.

\subsection{The smaller direct-coordinate quartic}
\label{int:direct}

For the preferred main quartic, retain the parent $J,w_{\rm scale}$ and
replace only $X$ by $3X_3-1$, while deleting $X-2u^2$ and $u$.
Denote this polynomial by $Q_{40}^{\mathrm{dir}}$; it is the complete
packet in \texttt{quartic40\_direct.json}.
Equivalently, use all twenty-five rows of Subsection~\ref{int:literal}
with the core list modified by $J_3\mapsto J$, $w_3\mapsto w_{\rm scale}$,
and $X\mapsto X_3$, and now use the abbreviations
\[
 J_0=J,\qquad q=\beta J+1,\qquad W=4w_{\rm scale},
 \qquad X=3X_3-1.
\]
The other abbreviations and the nine private guard coordinates are
unchanged.  The resulting sum of squares has $40$ unknowns, $25$
quadratic residuals, exact degree four, and $393$ expanded symbolic
monomials.  Its coefficient of $\tau^4$ remains one.

Here $X\equiv2\pmod3$ gives $X\neq0$ immediately.  The retained equation
$X=qW$ and $q\neq0$ then give $W\neq0$, so $|W|\geq4$ and
$|X|\geq4|q|$.  Thus every bound and every step of the sign-uniform
argument applies without change, recovering $X=e2^R$ with $R$ odd.
The residue selects $e=1$, after which
$u=2^{(R-1)/2}$ and $X_{\rm old}=3X_3-1$ restore the parent.
For completeness, canonical $X_{\rm old}=2^R$ with $R$ odd gives the
integer $X_3=(2^R+1)/3$; keep $J,w_{\rm scale}$ and all remaining
coordinates and discard $u$.  This direct version needs no separate
residue restriction on $J,w_{\rm scale}$ or on $\beta$.

The two-affine-coordinate version is retained for the projections below:
there $X=Wq$ is a monic polynomial definition.  Directly solving
$3X_3-1=Wq$ would require division by $3$, which is not an allowed
integer-polynomial restoration and is not used here.

\subsection{The exact sextic and octic projections}
\label{int:projections}

Deleting a row $v-P(\text{other coordinates})$ and substituting its
integer-polynomial right-hand side everywhere gives a bijection of integer
zero sets.  The inverse restores $v$.  The transformations below use the
centered-root and triangular elimination ideas of source62; their purpose
here is to transfer the new deletion to additional degree slices.

\paragraph{Thirty-four integers, degree six.}
Make the triangular change
$\tau_{\rm old}=\tau_{\rm tail}+E(c-\eta)$ and eliminate
\[
 \begin{gathered}
 P=\eta+\zeta+E(c-\eta),\quad N=q(L+1),\quad X=Wq,\\
 C_2=c^2,\qquad F_2=f^2,\qquad g=y^2-V^2.
 \end{gathered}
\]
The first norm becomes exactly
\[
 \tau_{\rm tail}^2+
 (2\tau_{\rm tail}-\eta-\zeta)E(c-\eta)-1.
\]
The leading quartic terms cancel, so every surviving row has degree at
most three.  Six coordinates and six defining rows disappear, leaving
thirty-four integers and nineteen rows.  The only row involving $r$ is
$I-(2r+1)c^2$, so the coefficient of $r^2c^4$ in their sum of squares
is exactly $4$.  The equation in \texttt{sextic34.json} has exact degree six.

\paragraph{Twenty-nine integers, degree eight.}
Put $k_0=1-V+jc+hE$ and eliminate the eleven coordinates
\[
\begin{gathered}
 L=q^2-1,\quad N=q^3,\quad X=Wq,\quad C_2=c^2,\quad F_2=f^2,\\
 g=y^2-V^2,\quad B_0=I^2,\quad
 D=Wq+ac+\gamma(4a+3),\\
 \Delta=a^2+4a+3,\quad
 \zeta=k_0-\eta,\quad P=k_0+E(c-\eta).
\end{gathered}
\]
Every right-hand side uses retained coordinates only.  The index-exposure
row $r_{17}$ supplies the monic definition of $\zeta$, while $r_7$
supplies the restoration of $P$.  Fourteen rows remain, all of degree at
most four.  The ordinary strong row
\[
 I^2-(a^2+4a+3)(f^2-1)
\]
is the only row containing both $a$ and $f$; its square contributes
$a^4f^4$ with coefficient $1$.  Therefore the equation in
\texttt{octic29.json} has exact degree eight.  Both projections preserve
the exact fixed tuple and inherit the same consistency equivalence.

\subsection{A general principle and what the comparison means}
\label{int:principle}

\begin{proposition}[Orientation after arithmetic recovery]
Suppose a polynomial system, without a designated sign guard, has an
arithmetic proof that every integer zero satisfies $X=\pm b^r$, where
$r$ lies in a specified set of integer indices.  Let $m$ be a fixed modulus
such that the possible residues of $b^r$ and $-b^r$ are disjoint.  If
affine restrictions on existing coordinates force $X$ into the positive
residue class, preserve a complete canonical witness family, and the
deleted guard can be restored from the recovered positive power, then
they replace that guard without introducing new unknowns.
\end{proposition}
\begin{proof}
Disjointness excludes the negative branch of the recovered power.
The original sign guard can then be restored if it follows from the
recovered positive power, as $X=2u^2$ does here for odd $r$.
Canonical completeness supplies the converse existence implication.
\end{proof}

In this construction $b=2$, $m=3$, and the index is odd: the two residue
sets are the singletons $\{2\}$ and $\{1\}$.  The important requirement
is the sign-uniform recovery theorem; simply dropping a square guard
and imposing a residue would be unjustified without it.

At the inspected source pin the comparable variable counts are
$41,35,30$ at degrees $4,6,8$.  The new counts are $40,34,29$,
with respectively $25,19,14$ residuals.  Their exact expanded sums of
squares, with the four fixed ports left as coefficient symbols, have
$393,614,1238$ monomials when the direct quartic is selected; the
two-affine-coordinate quartic used for projections has $614$.
The prior $41$-variable quartic has $388$.
The affine restrictions enlarge these expanded
expressions.  Thus the gain is in unrestricted-integer variables and
residual counts at fixed degree, not a decrease in every size measure.
Native coefficient payload size is unchanged; no new bound for nine
positive or ten unrestricted integer unknowns is asserted.

The standard-library builder and a separately authored verifier both check
all literal affine substitutions, all twenty-five restoration identities
for each projection, exact symbolic expansions, liveness, degrees,
nonzero degree certificates, source hashes, and the actual coefficient
payloads.  Both pass in normal and optimized Python execution.  These
checks establish the finite algebraic claims.  The unbounded Pell,
native-compiler, and syntactic-refutation implications remain the stated
mathematical proof and its named dependencies.
